\subsection{The Koide cone and the canonical period discriminant}
\label{koide-s6-subsection-ambient-discriminant}

For \(r=(r_1,r_2,r_3)\in\mathbb R^3\), set
\begin{equation}
 Q_K(r)
 :=\frac{2}{3}(r_1+r_2+r_3)^2-(r_1^2+r_2^2+r_3^2)
 \label{koide-s6-eq-koide-form}
\end{equation}
and introduce the orthogonal coordinates
\begin{equation}
 \begin{pmatrix}h\\x\\y\end{pmatrix}
 =\mathcal F\begin{pmatrix}r_1\\r_2\\r_3\end{pmatrix},
 \qquad
 \mathcal F=
 \begin{pmatrix}
  1/\sqrt3&1/\sqrt3&1/\sqrt3\\
  1/\sqrt2&0&-1/\sqrt2\\
  1/\sqrt6&-2/\sqrt6&1/\sqrt6
 \end{pmatrix}.
 \label{koide-s6-eq-orthogonal-coordinates}
\end{equation}

We use exactly the canonical period convention
\begin{equation}
 Z=\begin{pmatrix}6\mu&\tau\\ \beta&\mu\end{pmatrix},
 \qquad \Pi=[Z\mid I_2].
 \label{koide-s6-eq-canonical-complex-period}
\end{equation}
Write
\[
 t=\operatorname{Im}\tau,\qquad
 m=\operatorname{Im}\mu,\qquad
 \varrho=\operatorname{Im}\beta.
\]
Relative to the period-lattice basis
\[
 (\widehat\gamma,\widehat u,\widehat w,\widehat\delta)
\]
and the interleaved real codomain basis
\[
 (\operatorname{Re}z_1,\operatorname{Im}z_1,
  \operatorname{Re}z_2,\operatorname{Im}z_2),
\]
the realification of \(\Pi\) is
\begin{equation}
 P(\tau,\mu,\beta)=
 \begin{pmatrix}
  6\operatorname{Re}\mu&\operatorname{Re}\tau&1&0\\
  6m&t&0&0\\
  \operatorname{Re}\beta&\operatorname{Re}\mu&0&1\\
  \varrho&m&0&0
 \end{pmatrix}.
 \label{koide-s6-eq-canonical-real-period}
\end{equation}
For comparison with the Lorentz cone, define the auxiliary symmetric
coefficient matrix
\begin{equation}
 S(t,m,\varrho):=
 \begin{pmatrix}t&\sqrt6\,m\\ \sqrt6\,m&\varrho\end{pmatrix}.
 \label{koide-s6-eq-symmetric-coefficient-matrix}
\end{equation}
Thus \(S\) represents the same coefficient discriminant, but it is not the
canonical \(4\times4\) real period matrix \(P\).

\begin{theorem}[Ambient coefficient-space discriminant isomorphism]
\label{koide-s6-thm-ambient-discriminant}
The matrix \(\mathcal F\) is orthogonal, and
\begin{equation}
 Q_K(r)=h^2-x^2-y^2.
 \label{koide-s6-eq-lorentz-identity}
\end{equation}
The real-linear map
\begin{equation}
 \mathcal K(h,x,y)
 =\bigl(\operatorname{Im}\tau,
         \operatorname{Im}\mu,
         \operatorname{Im}\beta\bigr)
 =\left(h+y,\frac{x}{\sqrt6},h-y\right)
 \label{koide-s6-eq-coefficient-map}
\end{equation}
is an isomorphism from \(\mathbb R^3\) onto the ambient real coefficient
space, with inverse
\begin{equation}
 \mathcal K^{-1}(t,m,\varrho)
 =\left(\frac{t+\varrho}{2},\sqrt6\,m,
        \frac{t-\varrho}{2}\right).
 \label{koide-s6-eq-coefficient-map-inverse}
\end{equation}
For every choice of the real parts of \(\tau,\mu,\beta\),
\begin{equation}
 \det P(\tau,\mu,\beta)
 =\det S(t,m,\varrho)
 =t\varrho-6m^2.
 \label{koide-s6-eq-canonical-discriminant}
\end{equation}
Consequently, if
\((t,m,\varrho)=\mathcal K(h,x,y)\), then
\begin{equation}
 \det P(\tau,\mu,\beta)
 =\det S\bigl(\mathcal K(h,x,y)\bigr)
 =Q_K(r).
 \label{koide-s6-eq-period-koide-identity}
\end{equation}
Moreover, for \(r\in\mathbb R_{\geq0}^3\setminus\{0\}\), the equation
\(Q_K(r)=0\) is equivalent to the assertion that
\begin{equation}
 S\bigl(\mathcal K(\mathcal F(r))\bigr)
 =\begin{pmatrix}h+y&x\\x&h-y\end{pmatrix}
 \label{koide-s6-eq-positive-sector-matrix}
\end{equation}
is positive semidefinite of rank one.
\end{theorem}

\begin{proof}
The rows of \(\mathcal F\) have norm one and pairwise scalar product zero,
so \(\mathcal F\mathcal F^{\mathsf T}=I_3\).  Hence
\[
 r_1^2+r_2^2+r_3^2=h^2+x^2+y^2,
 \qquad r_1+r_2+r_3=\sqrt3\,h.
\]
Substitution in \cref{koide-s6-eq-koide-form} proves
\cref{koide-s6-eq-lorentz-identity}.

Adding and subtracting the first and third entries of
\cref{koide-s6-eq-coefficient-map}, and multiplying the second by
\(\sqrt6\), proves \cref{koide-s6-eq-coefficient-map-inverse}.  Expanding
\cref{koide-s6-eq-canonical-real-period} along its last two columns gives
\[
 \det P
 =-\det\begin{pmatrix}6m&t\\\varrho&m\end{pmatrix}
 =t\varrho-6m^2.
\]
Equivalently, in the exact canonical convention,
\(\det P=-\det\operatorname{Im}Z\).  Taking the determinant in
\cref{koide-s6-eq-symmetric-coefficient-matrix} proves the other equality
in \cref{koide-s6-eq-canonical-discriminant}, while
\[
 (h+y)(h-y)-6\left(\frac{x}{\sqrt6}\right)^2
 =h^2-x^2-y^2
\]
proves \cref{koide-s6-eq-period-koide-identity}.

Finally, let \(r\in\mathbb R_{\geq0}^3\setminus\{0\}\).  Then
\(h=(r_1+r_2+r_3)/\sqrt3>0\).  If \(Q_K(r)=0\), the symmetric matrix in
\cref{koide-s6-eq-positive-sector-matrix} has determinant zero and trace
\(2h>0\); its eigenvalues are therefore \(0\) and \(2h\), so it is positive
semidefinite of rank one.  Conversely, rank one forces its determinant,
and hence \(Q_K(r)\), to vanish.
\end{proof}

\begin{corollary}[The Koide cone is not a family of regular tori]
\label{koide-s6-cor-ambient-only}
The identification in
\cref{koide-s6-thm-ambient-discriminant} is only a real-linear isomorphism
of ambient coefficient spaces together with their quadratic discriminants.
It neither supplies the real parts of the periods nor imposes the
holomorphic and equivariance conditions of the canonical period family.
At every point of the Koide cone, every matrix of the canonical form
\cref{koide-s6-eq-canonical-real-period} with coefficient triple
\cref{koide-s6-eq-coefficient-map} satisfies \(\det P=0\).  Its columns
therefore do not span \(\mathbb R^4\), do not generate a full lattice, and
do not define a regular compact complex two-torus.  Thus the ambient
discriminant identity is not a map from Koide triples to regular fibres of
the canonical \(S^6\) family.
\end{corollary}

\begin{corollary}[Separation from the actual canonical cusp]
\label{koide-s6-cor-actual-cusp}
In the notation
\[
 L_T=T+H_T,\qquad q_T=L_T-B_T,\qquad
 D_T=L_Tq_T+6m_T^2
\]
of the positive radial cusp, the already-proved determinant formula
\cref{cusp-spec-determinant} gives
\begin{equation}
 \det P_T=-D_T<0.
 \label{koide-s6-eq-actual-cusp-determinant}
\end{equation}
Moreover, with
\(H=\operatorname{Im}h(0)\) and \(B=\operatorname{Im}b(0)\),
\cref{cusp-spec-D-expansion} gives
\begin{equation}
 D_T=T^2+(2H-B)T+O(1),
 \qquad
 \det P_T\sim-T^2.
 \label{koide-s6-eq-actual-cusp-asymptotic}
\end{equation}
Thus every actual radial cusp fibre is regular and has nonzero oriented
period determinant, whereas every point in the Koide image has determinant
zero.
\end{corollary}

\begin{proof}
The first assertion is exactly \cref{cusp-spec-determinant}; the stated
asymptotic follows from \cref{cusp-spec-D-expansion}.  Since \(D_T>0\), the
matrix \(P_T\) is invertible and its columns form the full period lattice of
the corresponding regular fibre.
\end{proof}
